
\item
Pour toute extension @@M0@@ de @@M1@@, le noyau de @@M2@@ est un groupe de torsion.
\item Pour toute extension pure @@M3@@ de @@M4@@, l'homomorphisme @@M5@@ est injectif. (NB. Noter qu'on ne suppose pas nécessairement @@M6@@ infini).
\end{enumerate}

Dans le cas (i), on est ramené à prouver que tout élément @@M7@@ dans le noyau de @@M8@@, où @@M9@@ est une algèbre intègre de type fini sur @@M10@@, est un élément de torsion. Utilisant un @@M11@@-homomorphisme @@M12@@, où @@M13@@ est une @@M14@@-algèbre finie, on est ramené au cas où @@M15@@ est une extension finie de @@M16@@. Alors on a @@M17@@ en vertu de la propriété postulée plus haut pour l'homomorphisme norme. Cela prouve que @@M18@@ est un élément de torsion. On aura même @@M19@@ dans le cas où le pgcd des degrés des extensions finies @@M20@@ de @@M21@@ quotients de @@M22@@ est égal à 1. Or cela est le cas si le corps des fractions @@M23@@ de @@M24@@ est une extension pure de @@M25@@, comme on constate aisément (seul le cas où @@M26@@ est fini demandant une nouvelle vérification). Cette remarque donne alors la conclusion de (ii).

\textbf{7.18.4.} Supposons que pour tout ouvert affine non vide @@M27@@ de l'espace affine @@M28@@, si @@M29@@ est l'anneau de @@M30@@, l'homomorphisme @@M31@@ est surjectif (comparer IX 1.6). Alors, utilisant la commutation de @@M32@@ aux limites inductives filtrantes, on trouve que pour toute extension pure @@M33@@ de @@M34@@, @@M35@@ est surjectif. Si on suppose de plus @@M36@@ infini, ou qu'on peut définir des homomorphismes norme comme dans 7.18.3, on en conclut que cet homomorphisme est même bijectif.

\textbf{7.18.5.} À titre d'exemple, considérons un schéma @@M37@@ de type fini sur @@M38@@, et posons
@@M39@@
